% Generated draft: requires mathematical and source review.
\documentclass[11pt]{article}
\usepackage[margin=25mm]{geometry}
\usepackage{fontspec}
\setmainfont{Latin Modern Roman}
\usepackage{amsmath,amssymb,mathtools,unicode-math}
\setmathfont{Latin Modern Math}
\usepackage{enumitem,needspace,xurl,hyperref}
\setlength{\parindent}{0pt}
\setlength{\parskip}{5pt}
\setlength{\emergencystretch}{3em}
\newcommand{\N}{\mathbb{N}}
\newcommand{\Z}{\mathbb{Z}}
\newcommand{\Q}{\mathbb{Q}}
\newcommand{\R}{\mathbb{R}}
\newcommand{\C}{\mathbb{C}}
\newcommand{\F}{\mathbb{F}}
\newcommand{\A}{\mathbb{A}}
\newcommand{\PP}{\mathbb{P}}
\newcommand{\E}{\mathbb{E}}
\newcommand{\eps}{\varepsilon}
\newcommand{\dd}{\,\mathrm{d}}
\begin{document}
\section*{UnsolvedMath ID: 30000907}
\textbf{Source problem code:} catalogue entry
\subsection{Definitions and mathematical statement}
Let $N\models\mathrm{ZFC}$. For each cardinal $\lambda$ as computed in $N$, let $\operatorname{Col}(\omega,\lambda)^N$ denote the intended collapse forcing in $N$ which makes $\lambda$ countable; a standard presentation uses finite partial functions from $\omega$ to $\lambda$, ordered by reverse inclusion. The precise presentation is not given in the record. Consider the property
\[
\forall\lambda\in\operatorname{Card}^N\quad
\forall H\text{ generic over }N\text{ for }\operatorname{Col}(\omega,\lambda)^N,
\qquad N[H]\equiv N.
\]
Elementary equivalence here means agreement on every parameter-free first-order sentence in the language of set theory:
\[
N[H]\models\sigma\quad\Longleftrightarrow\quad N\models\sigma
\quad\text{for every such sentence }\sigma.
\]
It does not assert $N\prec N[H]$ or agreement for formulas with arbitrary parameters from $N$. All cardinals of $N$ are quantified, not just a chosen cardinal or the regular cardinals. In a suitable formal treatment, generic invariance is expressed by the corresponding schema of forcing assertions, rather than by an unprovided internal truth predicate.

The problem asks for the consistency strength of the existence of such a model $N$. The original statement further conjectures that this strength is ``approximately the same'' as that of the hypothesis used in work cited by the report to construct a model in which the club filter on $\omega_1$ is an ultrafilter. In a model where this filter is defined, the club filter consists of subsets of $\omega_1$ containing a closed unbounded set; the ultrafilter requirement says that for each subset $A\subseteq\omega_1$, either $A$ or its complement contains such a set.

The comparison hypothesis and the cited construction are not identified in the supplied record. Its target model's axioms, notably choice principles, must also be specified; it cannot simply be identified with the ZFC model $N$ above. Nor does ``approximately the same'' specify an exact equiconsistency assertion or a pair of particular upper and lower bounds. The intended metatheory, formalization of models and generic extensions, and precise comparison theory remain necessary data.

The supplied literature assessment checked 2026-08-21 labels the question open. It says that an exact comparison with the ultrafilter hypothesis has not been established and that later work gives bounds rather than a definitive identification; this is dated attributed triage.
\subsection{Short English statement}
Determine the consistency strength of a ZFC model whose parameter-free first-order theory is unchanged by every collapse $\operatorname{Col}(\omega,\lambda)$ for its cardinals. Compare that strength with the report's unspecified hypothesis for a club-filter ultrafilter construction, retaining the original conjecture's approximate wording.
\subsection{Sources}
\begin{itemize}
\item Local UnsolvedMath registry, permanent ID 30000907, source code catalogue entry; supplied complete record. Its corrected-verified note reports direct confirmation of the forcing-invariance property and the requested comparison.
\item Set Theory (2008), Oberwolfach Reports, DOI 10.4171/owr/2008/02, pp. 79--118: \url{https://doi.org/10.4171/owr/2008/02}.
\item Golshani--Mitchell: \url{https://arxiv.org/abs/1609.02633}, supplied as later work relevant to the dated assessment.
\end{itemize}
\end{document}
